\documentclass[11pt]{article}

% Packages
\usepackage[utf8]{inputenc}
\usepackage[T1]{fontenc}
\usepackage{amsmath, amssymb, amsthm}
\usepackage{geometry}
\usepackage{fancyhdr}
\usepackage{enumitem}
\usepackage{graphicx}

% Page geometry
\geometry{margin=1in}

% Header/Footer
\pagestyle{fancy}
\fancyhf{}
\lhead{CS556 Final Exam}
\rhead{John Rizzo}
\cfoot{\thepage}

% Custom commands for common math notation
\newcommand{\R}{\mathbb{R}}
\newcommand{\N}{\mathbb{N}}
\newcommand{\Z}{\mathbb{Z}}
\newcommand{\Q}{\mathbb{Q}}
\newcommand{\E}{\mathbb{E}}
\newcommand{\Var}{\text{Var}}
\newcommand{\Cov}{\text{Cov}}
\newcommand{\argmax}{\operatorname{argmax}}
\newcommand{\argmin}{\operatorname{argmin}}

% Theorem environments
\theoremstyle{definition}
\newtheorem{problem}{Problem}
\newtheorem*{solution}{Solution}

\begin{document}

\begin{center}
    {\LARGE \textbf{CS556 Final Exam}} \\[0.5em]
    {\large John Rizzo} \\[0.3em]
    {\large \today}
\end{center}

\vspace{1em}

%----------------------------------------------------------------------

\begin{problem}
(10 points) If possible, give a matrix $A$ which has $\begin{bmatrix}
    1 \\ 2 \\ 1
\end{bmatrix}$ as a basis for the column space and $\begin{bmatrix}
    0 \\ 3 \\ 2 \\ -1
\end{bmatrix}$ as a basis for the row space. If not possible, give your reason.
\end{problem}

\begin{solution}
    It is not possible to have such a matrix $A$. The rank of a matrix is equal to the dimension of its column space and also equal to the dimension of its row space. 
    Since the rank of $A$ (which is 1) cannot exceed the number of rows or columns, it is impossible for $A$ to have a row space spanned by a 4-dimensional vector while having a column space spanned by a single vector. So, the matrix $A$ cannot exist.
\end{solution}

%----------------------------------------------------------------------

\begin{problem}
(10 points) Let $A = \begin{bmatrix}
    1 & 4 \\
    2 & 3
\end{bmatrix}$. Find an invertible matrix $S$ that makes $S^{-1}AS$ a diagonal matrix.
\end{problem}

\begin{solution}
    $A = \begin{bmatrix}
        1 & 4 \\
        2 & 3
    \end{bmatrix}$

    \begin{align}
        det(A - \lambda I) &= det\begin{bmatrix}
            1 - \lambda & 4 \\
            2 & 3 - \lambda
        \end{bmatrix} \\
        &= (1 - \lambda)(3 - \lambda) - 8 \\
        &= \lambda^2 - 4\lambda - 5 \\
        &= (\lambda - 5)(\lambda + 1)
    \end{align}

    The eigenvalues are $\lambda_1 = 5$ and $\lambda_2 = -1$.

    For $\lambda_1 = 5$:
    \begin{align}
        (A - 5I)\mathbf{v} &= 0 \\
        \begin{bmatrix}
            -4 & 4 \\
            2 & -2
        \end{bmatrix} \begin{bmatrix}
            x_1 \\ x_2
        \end{bmatrix} &= 0 \\
        -4x_1 + 4x_2 &= 0 \\
        2x_1 - 2x_2 &= 0
    \end{align}
    
    This gives the eigenvector $\mathbf{v_1} = \begin{bmatrix}
        1 \\ 1
    \end{bmatrix}$.

    For $\lambda_2 = -1$:
    \begin{align}
        (A + I)\mathbf{v} &= 0 \\
        \begin{bmatrix}
            2 & 4 \\
            2 & 4
        \end{bmatrix} \begin{bmatrix}
            x_1 \\ x_2
        \end{bmatrix} &= 0 \\
        2x_1 + 4x_2 &= 0 \\
        2x_1 + 4x_2 &= 0
    \end{align}

    This gives the eigenvector $\mathbf{v_2} = \begin{bmatrix}
        -2 \\ 1
    \end{bmatrix}$.

    Thus, we can form the matrix $S$ using the eigenvectors as columns:
    \begin{align}
        S &= \begin{bmatrix}
            1 & -2 \\
            1 & 1
        \end{bmatrix}
    \end{align}
\end{solution}

%----------------------------------------------------------------------

\begin{problem}
(10 points) Each independent question refers to the matrix $A = \begin{bmatrix}
    4 & 1 \\
    d & -4 
\end{bmatrix}$. In each case, find the value of $d$ that makes the statement true (and show your work!).
\begin{enumerate}[label=(\alph*)]
    \item Given a value of $d$ such that $\begin{bmatrix}
        5 \\ 1
    \end{bmatrix}$ is an eigenvector of $A$. 

    We can find the corresponding eigenvalue $\lambda$:
    \begin{align}
        A\begin{bmatrix}
            5 \\ 1
        \end{bmatrix} &= \lambda \begin{bmatrix}
            5 \\ 1
        \end{bmatrix} \\
        \begin{bmatrix}
            4 & 1 \\
            d & -4 
        \end{bmatrix} \begin{bmatrix}
            5 \\ 1
        \end{bmatrix} &= \lambda \begin{bmatrix}
            5 \\ 1
        \end{bmatrix} \\
        \begin{bmatrix}
            21 \\ 5d - 4
        \end{bmatrix} &= \begin{bmatrix}
            5\lambda \\ \lambda
        \end{bmatrix}
    \end{align}
    From the first row, we have:
    \begin{align}
        21 &= 5\lambda \\
        \lambda &= \frac{21}{5}
    \end{align}
    From the second row, we have:
    \begin{align}
        5d - 4 &= \lambda \\
        5d - 4 &= \frac{21}{5} \\
        5d &= \frac{21}{5} + 4 \\
        5d &= \frac{21}{5} + \frac{20}{5} \\
        5d &= \frac{41}{5} \\
        d &= \frac{41}{25}
    \end{align}

    \item Given a value of $d$ such that 2 is one of the eigenvalues of $A$.
    \begin{align}
        det(A - \lambda I) &= 0 \\
        det\begin{bmatrix}
            4 - \lambda & 1 \\
            d & -4 - \lambda
        \end{bmatrix} &= 0 \\
        (4 - \lambda)(-4 - \lambda) - d &= 0 \\
        \lambda^2 - 16 - d &= 0
    \end{align}
    Substituting $\lambda = 2$:
    \begin{align}
        (2)^2 - 16 - d &= 0 \\
        4 - 16 - d &= 0 \\
        -12 - d &= 0 \\
        d &= -12
    \end{align}

    \item Given a value for $d$ such that $A$ is a nondiagonalizable matrix.
    \begin{align}
        det(A - \lambda I) &= 0 \\
        det\begin{bmatrix}
            4 - \lambda & 1 \\
            d & -4 - \lambda
        \end{bmatrix} &= 0 \\
        (4 - \lambda)(-4 - \lambda) - d &= 0 \\
        \lambda^2 - 16 - d &= 0
    \end{align}

\end {enumerate}
\end{problem}

\begin{solution}
\end{solution}

%----------------------------------------------------------------------

\begin{problem}
(10 points)  A rectangular storage container with an open top is to have a volume of $10m^3$.  
The length of its base is twice the width.  The material for the base costs \$10 per square meter.  
Material for the sides costs \$6 per square meter.  Find the cost of materials for the cheapest such container.
\end{problem}

\begin{solution}
    The Volume of the container is calculated by:
    \begin{align}
        V &= l \cdot w \cdot h \\
        10 &= (2w) \cdot w \cdot h \\
        10 &= 2w^2h \\
        h &= \frac{5}{w^2}
    \end{align}
    The cost of the container is the cost of the base plus the cost of the sides:
    \begin{align}
        C &= C_{base} + C_{sides} \\
        C_{base} &= 10 \cdot (l \cdot w) = 10 \cdot (2w \cdot w) = 20w^2 \\
        C_{sides} &= 6 \cdot (2(lh) + 2(wh)) = 6 \cdot (2(2wh) + 2(wh)) = 6 \cdot (4wh + 2wh) = 6 \cdot 6wh = 36wh \\
        C_{sides} &= 36w \cdot \frac{5}{w^2} = \frac{180}{w} \\
        C &= 20w^2 + \frac{180}{w}
    \end{align}
    To minimize the cost, we take the derivative of $C$ with respect to $w$ and set it to zero:
    \begin{align}
        \frac{dC}{dw} &= 40w - \frac{180}{w^2} \\
        0 &= 40w - \frac{180}{w^2} \\
        40w^3 &= 180 \\
        w^3 &= \frac{9}{2} \\
        w &= \sqrt[3]{\frac{9}{2}} \approx 1.65 \, m
    \end{align}
    \begin{align}
        h &= \frac{5}{w^2} = \frac{5}{(\sqrt[3]{\frac{9}{2}})^2} = \frac{5}{(\frac{9}{2})^{2/3}} \approx 1.83 \,m
    \end{align}
    The minimum cost: \\
    $C = 20w^2 + \frac{180}{w}$ \\
    $C = 20(1.65)^2 + \frac{180}{1.65} \approx \$163.54$ \\

    The cost of materials for the cheapest container is approximately \$163.54.
\end{solution}

%----------------------------------------------------------------------

\begin{problem}
(10 points) A social scientist decides to study in depth one individual to be randomly sampled from a given population.  
The social scientist is worried about the possibility that the chosen individual will refuse to cooperate.  
A previous student provides data on the probability of noncooperation for individuals in different income classes.
\begin{center}
    \begin{tabular}{l|c}
        High Income & 0.08 \\
        Middle Income & 0.10 \\
        Low Income & 0.04
    \end{tabular}
\end{center}
The scientist knows that in the population the proportion of high income people is 0.25, of medium income people is 0.5, 
and low income people is 0.25.
\begin{enumerate}
    \item What is the probability that a randomly selected individual will be noncooperative?
    \begin{align}
        P(NC) &= P(NC|HI)P(HI) + P(NC|MI)P(MI) + P(NC|LI)P(LI) \\
        &= (0.08)(0.25) + (0.10)(0.5) + (0.04)(0.25) \\
        &= 0.02 + 0.05 + 0.01 \\
        &= 0.08
    \end{align}
    \item Assuming that a randomly selected individual is noncooperative, what is the probability that that individual has low income?
    The probability that a randomly selected individual is noncooperative and has low income is given by:
    \begin{align}
        P(NC \cap LI) &= P(NC|LI)P(LI) \\
        &= (0.04)(0.25) \\
        &= 0.01
    \end{align}
    Using Bayes' theorem, we can find the probability that a randomly selected individual has low income given that they are noncooperative:
    \begin{align}
        P(LI|NC) &= \frac{P(NC \cap LI)}{P(NC)} \\
        &= \frac{0.01}{0.08} \\
        &= 0.125
    \end{align}
\end{enumerate}
\end{problem}

%----------------------------------------------------------------------

\begin{problem}
(10 points) If the function $f(x) = x^3 + ax^2 + bx$ has the local minimum value of $\frac{-2}{9}\sqrt{3}$ at $x = \frac{1}{\sqrt{3}}$, 
what are the values of $a$ and $b$.
\end{problem}

\begin{solution}
    \begin{align}
        f'(x) &= 3x^2 + 2ax + b
    \end{align}
    Since $f(x)$ has a local minimum at $x = \frac{1}{\sqrt{3}}$, we know that $f'(\frac{1}{\sqrt{3}}) = 0$:
    \begin{align}
        f'(\frac{1}{\sqrt{3}}) &= 3(\frac{1}{\sqrt{3}})^2 + 2a(\frac{1}{\sqrt{3}}) + b = 0 \\
        3(\frac{1}{3}) + \frac{2a}{\sqrt{3}} + b &= 0 \\
        1 + \frac{2a}{\sqrt{3}} + b &= 0 \\
        \frac{2a}{\sqrt{3}} + b &= -1 \quad (1)
    \end{align}
    Next, we know that $f(\frac{1}{\sqrt{3}}) = \frac{-2}{9}\sqrt{3}$:
    \begin{align}
        f(\frac{1}{\sqrt{3}}) &= (\frac{1}{\sqrt{3}})^3 + a(\frac{1}{\sqrt{3}})^2 + b(\frac{1}{\sqrt{3}}) = \frac{-2}{9}\sqrt{3} \\
        \frac{1}{3\sqrt{3}} + a(\frac{1}{3}) + b(\frac{1}{\sqrt{3}}) &= \frac{-2}{9}\sqrt{3} \\
        \frac{1}{3\sqrt{3}} + \frac{a}{3} + \frac{b}{\sqrt{3}} &= \frac{-2}{9}\sqrt{3} \\
        \frac{a}{3} + \frac{b}{\sqrt{3}} &= \frac{-2}{9}\sqrt{3} - \frac{1}{3\sqrt{3}} \\
        \frac{a}{3} + \frac{b}{\sqrt{3}} &= \frac{-6 - 1}{9\sqrt{3}} \\
        \frac{a}{3} + \frac{b}{\sqrt{3}} &= \frac{-7}{9\sqrt{3}} \quad (2)
    \end{align}
    Multiplying equation (1) by $\sqrt{3}$ gives:
    \begin{align}
        2a + b\sqrt{3} &= -\sqrt{3} \quad (3)
    \end{align}
    Multiplying equation (2) by $3\sqrt{3}$ gives:
    \begin{align}
        a\sqrt{3} + 3b &= -\frac{7}{3} \quad (4)
    \end{align}
    Multiplying equation (3) by 3 gives:
    \begin{align}
        6a + 3b\sqrt{3} &= -3\sqrt{3} \quad (5)
    \end{align}
    Multiplying equation (4) by $\sqrt{3}$ gives:
    \begin{align}
        a \cdot 3 + 3b\sqrt{3} &= -\frac{7\sqrt{3}}{3} \quad (6)
    \end{align}
    Subtracting equation (6) from equation (5) gives:
    \begin{align}
        3a &= -3\sqrt{3} + \frac{7\sqrt{3}}{3} \\
        3a &= \frac{-9\sqrt{3} + 7\sqrt{3}}{3} \\
        3a &= \frac{-2\sqrt{3}}{3} \\
        a &= \frac{-2\sqrt{3}}{9} \\
        a &= -2
    \end{align}
    Substituting $a = -2$ into equation (1) gives:
    \begin{align}
        \frac{2(-2)}{\sqrt{3}} + b &= -1 \\
        \frac{-4}{\sqrt{3}} + b &= -1 \\
        b &= -1 + \frac{4}{\sqrt{3}} \\
        b &= 1
    \end{align}
    The values of $a$ and $b$ are $-2$ and $1$ respectively.
\end{solution}

%----------------------------------------------------------------------

\begin{problem}
(10 points) Consider the function: $f(x,y) = x^2y + 3sin(xy) + 2y^2$.  We apply gradient descent with a fixed learning rate $\alpha = 0.1$
starting from the initial point $(x_0, y_0) = (1,1)$. Using gradient descent, compute the points $(x_1, y_1)$, $(x_2, y_2)$, and $(x_3, y_3)$
reached after three iterations.
\end{problem}

\begin{solution}
    The gradient of $f(x,y)$:
    \begin{align}
        \frac{\partial f}{\partial x} &= 2xy + 3ycos(xy) \\
        \frac{\partial f}{\partial y} &= x^2 + 3xcos(xy) + 4y
    \end{align}
    The gradient at the initial point $(x_0, y_0) = (1,1)$:
    \begin{align}
        \nabla f(1,1) &= \begin{bmatrix}
            2(1)(1) + 3(1)cos(1) \\
            (1)^2 + 3(1)cos(1) + 4(1)
        \end{bmatrix} \\
        &= \begin{bmatrix}
            2 + 3cos(1) \\
            1 + 3cos(1) + 4
        \end{bmatrix} \\
        &= \begin{bmatrix}
            2 + 3cos(1) \\
            5 + 3cos(1)
        \end{bmatrix}
    \end{align}
    Perform the gradient descent updates:
    \begin{align}
        (x_1, y_1) &= (x_0, y_0) - \alpha \nabla f(x_0, y_0) \\
        &= (1, 1) - 0.1 \begin{bmatrix}
            2 + 3cos(1) \\
            5 + 3cos(1)
        \end{bmatrix} \\
        &= (1 - 0.1(2 + 3cos(1)), 1 - 0.1(5 + 3cos(1))) \\
        &= (1 - 0.2 - 0.3cos(1), 1 - 0.5 - 0.3cos(1)) \\
        &= (0.8 - 0.3cos(1), 0.5 - 0.3cos(1))
    \end{align}
    \begin{align}
        \nabla f(x_1, y_1) &= \begin{bmatrix}
            2x_1y_1 + 3y_1cos(x_1y_1) \\
            x_1^2 + 3x_1cos(x_1y_1) + 4y_1
        \end{bmatrix} \\
        (x_2, y_2) &= (x_1, y_1) - 0.1 \nabla f(x_1, y_1) \\
        \nabla f(x_2, y_2) &= \begin{bmatrix}
            2x_2y_2 + 3y_2cos(x_2y_2) \\
            x_2^2 + 3x_2cos(x_2y_2) + 4y_2
        \end{bmatrix} \\
        (x_3, y_3) &= (x_2, y_2) - 0.1 \nabla f(x_2, y_2)
    \end{align}
\end{solution}

%----------------------------------------------------------------------

\begin{problem}
(5 points)  Find the derivative of $(4x^2 - 7)^{2+\sqrt{x^2 - 5}}$ using logarithmic differentiation.
\end{problem}

\begin{solution}
    \begin{align}
        log(y) &= log((4x^2 - 7)^{2+\sqrt{x^2 - 5}}) = (2+\sqrt{x^2 - 5})log(4x^2 - 7) \\
        \frac{1}{y} \frac{dy}{dx} &= \frac{d}{dx}[(2+\sqrt{x^2 - 5})log(4x^2 - 7)] \\
        \frac{1}{y} \frac{dy}{dx} &= log(4x^2 - 7) \frac{d}{dx}(2+\sqrt{x^2 - 5}) + (2+\sqrt{x^2 - 5}) \frac{d}{dx}[log(4x^2 - 7)] \\
        \frac{1}{y} \frac{dy}{dx} &= log(4x^2 - 7) \frac{1}{2\sqrt{x^2 - 5}}(2x) + (2+\sqrt{x^2 - 5}) \frac{1}{4x^2 - 7}(8x) \\
        \frac{dy}{dx} &= y [ log(4x^2 - 7) \frac{x}{\sqrt{x^2 - 5}} + (2+\sqrt{x^2 - 5}) \frac{8x}{4x^2 - 7} ] \\
        \frac{dy}{dx} &= (4x^2 - 7)^{2+\sqrt{x^2 - 5}} [ log(4x^2 - 7) \frac{x}{\sqrt{x^2 - 5}} + (2+\sqrt{x^2 - 5}) \frac{8x}{4x^2 - 7} ]
    \end{align}
\end{solution}

%----------------------------------------------------------------------

\begin{problem}
(10 points) The breakdown voltage of a randomly chosen diode of a certain type is known to be normally distributed with a mean value of 
40 volt and standard deviation of 1.5 volt.
\begin{enumerate}
    \item What is the probability that the voltage of a single diode is between 39 and 42?
    \item What value is such that only 15\% of all diodes have voltages exceedign that value?
    \item If four diodes are independently selected, what is the probability that at least one has a voltage exceeding 42?
\end{enumerate}
\end{problem}

\begin{solution}
    The proability that the voltage of a single diode is between 39 and 42 is approximately 0.5793.
    \begin{align}
        Z_1 = \frac{39 - 40}{1.5} &\approx -0.6667 \\
        Z_2 = \frac{42 - 40}{1.5} &\approx 1.3333 \\
        P(Z < -0.6667) &\approx 0.2525 \\
        P(Z < 1.3333) &\approx 0.9082
    \end{align}
    The probability that the voltage is between 39 and 42 is:
    \begin{align}
        P(39 < X < 42) &= P(Z < 1.3333) - P(Z < -0.6667) \\
        &\approx 0.9082 - 0.2525 \\
        &\approx 0.6557
    \end{align}
\end{solution}

\begin{solution}
    \item The value such that only 15\% of all diodes have voltages exceeding that value is approximately 42.06 volts.
    We need to find the z-score corresponding to the 85th percentile (since 15\% exceed that value):
    \begin{align}
        Z_{0.85} &\approx 1.036 \\
        X &= \mu + Z\sigma \\
        X &= 40 + 1.036 \cdot 1.5 \\
        X &\approx 40 + 1.554 \\
        X &\approx 41.554
    \end{align}
\end{solution}
\begin{solution}
    \item The probability that a single diode has a voltage exceeding 42:
    \begin{align}
        Z &= \frac{42 - 40}{1.5} \approx 1.3333
    \end{align}
    \begin{align}
        P(Z < 1.3333) &\approx 0.9082 \\
        P(X > 42) &= 1 - P(Z < 1.3333) \\
        &\approx 1 - 0.9082 \\
        &\approx 0.0918
    \end{align}
    \begin{align}
        P(\text{none exceed 42}) &= (1 - P(X > 42))^4 \\
        &\approx (1 - 0.0918)^4 \\
        &\approx (0.9082)^4 \\
        &\approx 0.6811
    \end{align}
    \begin{align}
        P(\text{at least one exceeds 42}) &= 1 - P(\text{none exceed 42}) \\
        &\approx 1 - 0.6811 \\
        &\approx 0.3189
    \end{align}
\end{solution}
%----------------------------------------------------------------------

\begin{problem}
(15 points) Find the values of $a,b$ and $c$ if the parabola $y=ax^2+bx+c$, is to be tangent to the line $y=4x+1$, at the point (-1, -3),
and is to have a critical point at x = -2.
\end{problem}

\begin{solution}
    Since the parabola is tangent to the line at the point (-1, -3), we know that:
    \begin{align}
        ax^2 + bx + c &= 4x + 1 \quad \text{at } x = -1 \\
        a(-1)^2 + b(-1) + c &= 4(-1) + 1 \\
        a - b + c &= -4 + 1 \\
        a - b + c &= -3 \quad (1)
    \end{align}
    Since the parabola has a critical point at x = -2, we know that the derivative of the parabola is zero at that point:
    \begin{align}
        \frac{dy}{dx} &= 2ax + b \\
        2a(-2) + b &= 0 \\
        -4a + b &= 0 \\
        b &= 4a \quad (2)
    \end{align}
    \begin{align}
        a - (4a) + c &= -3 \\
        -3a + c &= -3 \\
        c &= 3a - 3 \quad (3)
    \end{align}
    To solve for $a$. The slope of the tangent line at the point (-1, -3) must equal the derivative of the parabola at that point:
    \begin{align}
        \frac{dy}{dx} &= 2ax + b \\
        2a(-1) + b &= 4 \\
        -2a + b &= 4 \quad (4)
    \end{align}
    \begin{align}
        -2a + (4a) &= 4 \\
        2a &= 4 \\
        a &= 2
    \end{align}
    Substitute $a = 2$ into equations (2) and (3) to find $b$ and $c$:
    \begin{align}
        b &= 4(2) = 8 \\
        c &= 3(2) - 3 = 6 - 3 = 3
    \end{align}
    The values of $a$, $b$, and $c$ are $2$, $8$, and $3$ respectively.
\end{solution}

%----------------------------------------------------------------------

\begin{problem}
(0 points) Would you like to be considered for a Course Assistant position for this course in Spring 2026?
\begin{enumerate}
    \item Yes
    \item No
\end{enumerate}
\end{problem}

\begin{solution}
\end{solution}
I would consider it an honor to be considered for a Course Assistant position for CS556 in Spring 2026.  I know that I have a 
lot of work to do in order to improve my mastery of the material, but I am committed to doing so.  I believe that being a Course Assistant
would help me to achieve that goal, and I would be grateful for the opportunity.  In addition I have substantial experience creating material
and training programs to introduce new technologies to software developers and other technical professionals in my work outside of school, and I believe that
this experience would help me to be an effective Course Assistant.
\end{document}
